\section{绪论}
\subsection{研究问题与研究边界}
旅游服务研究宜从具体场景提出问题，明确服务环节、参与者和分析单位。本例以服务信息传递为写作线索，演示文经管法类编号和研究论文结构，不报告真实问卷、访谈或企业资料。城市、机构和调查人数应在研究完成后如实填写。
\subsection{文献综述与概念界定}
综述需要比较概念、方法、证据及局限，不能把一般经验直接称为理论。资料应追溯至原始出版物，记录版本、页码和允许用途。本例使用真实出版物演示引文排版\cite{knuth1984,lamport1994}，不把它们冒充为旅游服务研究综述。
\subsubsection{研究对象与分析单位}
明确分析对象是游客个体、单次服务事件、团队还是机构。同一游客多次反馈不自动等于多个独立样本；便利样本也不代表所有游客。涉及个人资料时，应取得适当许可并匿名化\footnote{本例没有接触真实参与者。脚注只演示同页补充说明与每页重置编号。}。
\paragraph{四级标题示例}
本类编号采用“一、㈠、1、⑴”，目录仅列到三级。不要因模板提供四级编号就把每个自然段都设为标题。
\section{研究设计}
\subsection{资料来源与抽样}
报告资料来源、时间范围、选择方式和排除规则。访谈提纲应经过试访修订，问卷指标应说明来源及适用条件，不能随意拼接后声称具有经过验证的信效度。
\subsection{分析维度与编码规则}
表~\ref{tab:tourism} 是观察维度表的结构示例，并非实际调查结论。正式研究需给出类别边界、难以判断的情况及复核过程。
\begin{table}[htbp]\centering\caption{服务信息研究维度示例}\label{tab:tourism}
\begin{tabularx}{\linewidth}{lXX}\toprule 维度 & 记录内容 & 解释边界\\\midrule
信息获取 & 来源、时间点与可访问性 & 区分公开材料与个人体验\\
服务交互 & 具体事件与沟通过程 & 不以单例推断总体水平\\
问题处理 & 反馈渠道与响应过程 & 记录上下文与反例\\\bottomrule\end{tabularx}
\end{table}
\subsection{研究流程}
图~\ref{fig:tourism} 展示从问题定义到结果解释的组织方式。每一步均应留下过程材料，以便核查。
\begin{figure}[htbp]\centering
\begin{tikzpicture}[node distance=6mm,b/.style={draw,minimum width=65mm,minimum height=9mm,font=\small},>=Stealth]
\node[b](a){明确问题与取得必要许可};\node[b,below=of a](b){试用工具与收集资料};
\node[b,below=of b](c){整理、匿名化与编码复核};\node[b,below=of c](d){解释证据与讨论局限};
\draw[->](a)--(b);\draw[->](b)--(c);\draw[->](c)--(d);
\end{tikzpicture}\caption{研究流程示例}\label{fig:tourism}
\end{figure}
\subsection{统计口径与质性分析}
采用比例时，应说明分母、缺失与多选情况；采用质性分析时，应呈现主题形成、反例和解释路径。若某类事件数为 $n_k$，满足纳入条件的总事件数为 $N$，其比例示例如下：
\begin{equation}p_k=\frac{n_k}{N}\times100\%.\end{equation}
没有实际数据时不填入百分比，更不能虚构统计显著性。统计指标只是工具，不能替代研究问题和语境分析。
\section{结果报告}
\subsection{描述性结果}
正式报告应先说明资料数量、覆盖范围及排除情况，再按研究问题呈现发现。摘录应匿名并保留内部编号，不应断章取义。本示例没有真实调查，因此只提供报告结构，没有“发现的问题”可供引用。
\subsection{证据、观点与解释}
区分研究者直接观察、参与者陈述和研究者解释。负面反馈不等于所有服务均有问题，满意评价也不能证明某项管理措施产生因果效果。图表与文字应互相补充，不重复堆放同一信息。
\section{讨论与建议}
\subsection{与已有研究的关系}
比较研究结果与既有研究的一致或差异，说明场景、抽样和测量方法可能带来的影响。若提出服务改进建议，应明确实施主体、成本、风险及可检验目标，避免空泛对策。
\subsection{局限与后续工作}
说明样本选择、研究者参与、资料缺失和隐私限制等问题。未来计划不能写成已完成的成果。本工作稿不支持任何关于旅游服务质量的实证判断。
